\chapter{Funkcje: jednoznaczność przyporządkowania}\label{ch:functions}
\section{Od relacji do funkcji}
Relacja może przyporządkować jednemu argumentowi wiele wyników albo żadnego. Gdy podajemy wzór $x\mapsto x^2$, oczekujemy jednak, że dla każdego dopuszczalnego argumentu istnieje \emph{dokładnie jeden} wynik. To właśnie odróżnia funkcję od dowolnej relacji.
\begin{definicja}[Funkcja między zbiorami]\label{def:function}
Niech $A,B$ będą zbiorami. Aby określić \emph{odwzorowanie}
$f:A\to B$, podajemy jego dziedzinę $A$, przeciwdziedzinę $B$
i \emph{wykres} $G_f\subseteq A\times B$, spełniający warunek
\[
\forall a\in A\ \exists!b\in B\ ((a,b)\in G_f).
\]
Równoważnie, dane odwzorowanie można formalnie traktować jako
trójkę $(A,B,G_f)$. W dalszej części tekstu dla wygody
mówimy po prostu „funkcja $f:A\to B$” i operujemy jej wykresem,
ale pamiętamy o zadanej przeciwdziedzinie.
Piszemy $f(a)=b$ wtedy i tylko wtedy, gdy $(a,b)\in G_f$.
\emph{Zbiorem wartości} jest $f[A]:=\{f(a):a\in A\}\subseteq B$.
Symbol $B^A$ oznacza zbiór wszystkich funkcji z $A$ do $B$.
\end{definicja}

\begin{uwaga}[Istnienie zbioru wszystkich funkcji]\label{rem:function-set-zf-xliv}
Zapis $B^A$ nie oznacza jedynie nieformalnej kolekcji. Ponieważ
iloczyn $A\times B$ jest zbiorem (zob. konstrukcję w
Rozdziale~\ref{ch:relations}), można utworzyć zbiór
$\mathcal P(A\times B)$ jego podzbiorów. Następnie przez wyróżnianie
otrzymujemy
\[
 B^A=\{G\in\mathcal P(A\times B):
 \forall a\in A\ \exists!b\in B\ ((a,b)\in G)\}.
\]
Jest to zapis zbioru \emph{wykresów}. Jeśli odwzorowania formalizujemy
jako trójki $(A,B,G)$, to ustalone $A,B$ pozwalają jednoznacznie
odtworzyć trójkę z wykresu; dlatego oba ujęcia dają ten sam zbiór funkcji
po naturalnym utożsamieniu. Nie stosujemy tutaj aksjomatu wyboru.
\end{uwaga}

\begin{twierdzenie}[Kiedy istnieje funkcja między zadanymi zbiorami?]
\label{thm:function-existence-xliv}
Niech $A$ i $B$ będą zbiorami. Funkcja $A\to B$ istnieje wtedy
i tylko wtedy, gdy $A=\varnothing$ lub $B\ne\varnothing$.
Gdy $A=\varnothing$, istnieje dokładnie jedna taka funkcja, niezależnie
od $B$.
\end{twierdzenie}
\begin{proof}
Jeśli $A=\varnothing$, relacja pusta jest podzbiorem
$A\times B=\varnothing$. Warunek istnienia dokładnie jednej wartości
dla każdego argumentu jest spełniony, ponieważ w $A$ nie ma
argumentów. Innej relacji zawartej w $A\times B$ nie ma.
Jeśli $A\ne\varnothing$ i $f:A\to B$ istnieje, weźmy $a_0\in A$.
Wtedy $f(a_0)\in B$, a więc $B\ne\varnothing$.
Odwrotnie, jeśli $B\ne\varnothing$, ustalmy jeden element $b_0\in B$.
Przyporządkowanie $a\mapsto b_0$ dla wszystkich $a\in A$ jest funkcją.
Wybór jednego elementu z jednego niepustego zbioru nie wymaga aksjomatu
wyboru.
\end{proof}

\begin{uwaga}
Przeciwdziedzina jest częścią danych odwzorowania, ale \emph{nie musi
wynikać z samego wykresu}. Ten sam wzór $f(x)=x^2$ i ten sam wykres
mogą być rozważane raz jako odwzorowanie $\R\to\R$, a raz jako
odwzorowanie $\R\to[0,\infty)$; tylko drugie jest surjekcją.
Jeśli ustalimy tę samą dziedzinę i przeciwdziedzinę, dwa odwzorowania
$f,g:A\to B$ są równe wtedy i tylko wtedy, gdy $f(a)=g(a)$
dla każdego $a\in A$. Jest to równoznaczne z równością ich wykresów.
W konwencji trójek odwzorowania o różnych przeciwdziedzinach
są różnymi danymi, nawet gdy mają identyczny wykres.
\end{uwaga}
\begin{przyklad}[Relacja niebędąca funkcją]
Niech $A=\{1,2\}$, $B=\{u,v\}$. Relacja $\{(1,u),(1,v),(2,u)\}$ nie jest funkcją $A\to B$, ponieważ argumentowi $1$ przypisuje dwa różne wyniki. Relacja $\{(1,u)\}$ również nie jest funkcją $A\to B$: argument $2$ nie ma obrazu. Natomiast $\{(1,u),(2,u)\}$ jest funkcją, mimo że oba argumenty mają ten sam obraz.
\end{przyklad}

\begin{uwaga}[Trzy obiekty, których nie należy mylić]
Dla $f:A\to B$ rozróżniamy: argument $a\in A$,
jego pojedynczą wartość $f(a)\in B$ oraz zbiór wszystkich
wartości $f[A]\subseteq B$. Na przykład dla $f(x)=x^2$
na $A=\{-2,-1,0,1,2\}$ mamy $f(-2)=4$ oraz
$f[A]=\{0,1,4\}$. W szczególności zbiór wartości może mieć
mniej elementów niż dziedzina i nie musi być równy przeciwdziedzinie.
\end{uwaga}

\begin{uwaga}[Zależność iniekcji i surjekcji od danych]
Iniekcyjność zależy tylko od dziedziny i przypisanych wartości,
podczas gdy surjekcyjność zależy również od przeciwdziedziny.
Zmniejszenie przeciwdziedziny do $f[A]$ zawsze daje surjekcję
$A\to f[A]$. Nie wolno jednak w ten sposób automatycznie
zmieniać funkcji bez odnotowania zmiany jej danych.
\end{uwaga}



\begin{twierdzenie}[Liczba funkcji między zbiorami skończonymi]
\label{thm:finite-function-count-xliv}
Niech $A,B$ będą zbiorami skończonymi, $n,m\in\N$ oraz
$|A|=n$, $|B|=m$. Wtedy liczba funkcji $A\to B$ wynosi $m^n$. Przyjmujemy tu konwencję
$0^0=1$, wyrażającą istnienie jednej funkcji $\varnothing\to\varnothing$.
\end{twierdzenie}
\begin{proof}
Dla $n=0$ zbiór $A$ jest pusty. Z
Twierdzenia~\ref{thm:function-existence-xliv} istnieje dokładnie jedna
funkcja $A\to B$, zgodnie z wartością $m^0=1$, również dla $m=0$.
Załóżmy, że wzór wykazano dla dziedzin $n$-elementowych.
Niech $|A|=n+1$. Wybierzmy $a_0\in A$ i oznaczmy
$A'=A\setminus\{a_0\}$, wobec czego $|A'|=n$.
Funkcję $f:A\to B$ wyznaczają dokładnie dwie dane:
jej ograniczenie $f|_{A'}:A'\to B$ oraz wartość $f(a_0)\in B$.
Odwrotnie, dowolna taka para danych określa jedną funkcję na $A$,
ponieważ $A'$ i $\{a_0\}$ są rozłączne. Istnieje $m^n$ możliwości
ograniczenia i $m$ możliwości wartości w $a_0$; razem daje to
$m^n\cdot m=m^{n+1}$. Jest to rozumowanie skończone i nie
odwołuje się do aksjomatu wyboru.
\end{proof}
\begin{przyklad}[Ręczne liczenie funkcji]
Dla $A=\{1,2\}$ i $B=\{u,v,w\}$ funkcję określa para
$(f(1),f(2))$. Możliwe pary to
$(u,u),(u,v),(u,w),(v,u),(v,v),(v,w),(w,u),(w,v),(w,w)$.
Jest ich $9=3^2$. Gdy $B=\varnothing$, nie ma żadnej funkcji
z tego niepustego $A$; gdy $A=\varnothing$, istnieje jedna
funkcja nawet dla $B=\varnothing$.
\end{przyklad}

\section{Iniekcje, surjekcje i bijekcje}
\begin{definicja}[Trzy podstawowe własności]\label{def:inj-surj-bij}
Niech $f:A\to B$ będzie funkcją. Mówimy, że $f$ jest \emph{iniekcją}, jeśli $\forall a,a'\in A\ (f(a)=f(a')\Rightarrow a=a')$. Mówimy, że jest \emph{surjekcją}, jeśli $\forall b\in B\ \exists a\in A\ (f(a)=b)$. Jest \emph{bijekcją}, jeśli jest jednocześnie iniekcją i surjekcją.
\end{definicja}
\begin{przyklad}
Funkcja $f:\Z\to\Z$, $f(n)=2n$, jest iniekcją: z $2m=2n$ wynika $m=n$. Nie jest surjekcją, gdyż liczba $1$ nie ma przeciwobrazu. Funkcja $g:\Z\to\{0,1\}$, przypisująca resztę z dzielenia przez $2$, jest surjekcją, lecz nie jest iniekcją, bo $g(0)=g(2)$. Funkcja $h:\Z\to\Z$, $h(n)=n+1$, jest bijekcją: różne argumenty mają różne obrazy, a dla każdego $k\in\Z$ zachodzi $h(k-1)=k$.
\end{przyklad}

\begin{przyklad}[Sprawdzamy warunki po kolei]
Niech $A=\{1,2,3\}$, $B=\{a,b\}$ i
$f(1)=a$, $f(2)=a$, $f(3)=b$.
Aby sprawdzić iniekcję, porównujemy wartości dla różnych
argumentów: $f(1)=f(2)$, mimo że $1\ne2$; funkcja nie
jest więc iniekcją. Aby sprawdzić surjekcję, bierzemy
po kolei elementy przeciwdziedziny: $a=f(1)$ i $b=f(3)$,
więc każdy z nich jest osiągany i funkcja jest surjekcją.
Jeśli zmienimy przeciwdziedzinę na $\{a,b,c\}$, zachowując
te same wartości, funkcja przestanie być surjekcją, bo $c$
nie ma przeciwobrazu. Nie zmieniły się wzory, zmieniły się
dane określające funkcję.
\end{przyklad}

\begin{twierdzenie}[Charakteryzacja iniekcji przez włókna]\label{thm:fibres}
Niech $A,B$ będą zbiorami, a $f:A\to B$ funkcją.
Dla $b\in B$ zdefiniujmy włókno $f^{-1}[\{b\}]:=\{a\in A:f(a)=b\}$. Funkcja $f$ jest iniekcją wtedy i tylko wtedy, gdy każde włókno ma co najwyżej jeden element. Jest surjekcją wtedy i tylko wtedy, gdy każde włókno jest niepuste.
\end{twierdzenie}
\begin{proof}
Jeśli $f$ jest iniekcją i $a,a'$ należą do tego samego włókna, to $f(a)=b=f(a')$, a więc $a=a'$. Odwrotnie, gdy $f(a)=f(a')=b$, oba argumenty leżą we włóknie nad $b$; założenie o co najwyżej jednym elemencie daje $a=a'$. Druga równoważność wynika bezpośrednio z rozwinięcia definicji surjekcji: $b$ należy do zbioru wartości dokładnie wtedy, gdy istnieje $a$ z $f(a)=b$.
\end{proof}

\begin{przyklad}[Włókna w poprzednim przykładzie]
Dla funkcji z poprzedniego przykładu
$f^{-1}[\{a\}]=\{1,2\}$ oraz $f^{-1}[\{b\}]=\{3\}$.
Włókno nad $a$ ma dwa elementy, więc funkcja nie jest iniekcją.
Obydwa włókna nad elementami $B=\{a,b\}$ są niepuste,
więc funkcja jest surjekcją. W ten sposób twierdzenie o włóknach
pozwala rozpoznać obie własności bez zgadywania.
\end{przyklad}


\begin{twierdzenie}[Włókna rozcinają dziedzinę]\label{thm:fiber-partition-xliv}
Niech $f:A\to B$ będzie funkcją i niech $C=f[A]$. Rodzina
$\{f^{-1}[\{b\}]:b\in C\}$ jest podziałem $A$: każde jej włókno
jest niepuste, dwa różne włókna są rozłączne, a ich suma jest równa $A$.
Dla $A=\varnothing$ otrzymujemy pustą rodzinę, która jest podziałem
pustego zbioru w rozumieniu Definicji~\ref{def:partition}.
\end{twierdzenie}
\begin{proof}
(1) Jeśli $b\in C$, to z definicji obrazu istnieje $a\in A$
taki, że $f(a)=b$. Stąd $a\in f^{-1}[\{b\}]$ i włókno jest niepuste.

(2) Jeśli $b,c\in C$ i $b\ne c$, to żaden $a\in A$ nie może
należeć do obu włókien: dawałoby to jednocześnie $f(a)=b$ i $f(a)=c$,
wbrew $b\ne c$.

(3) Każdy $a\in A$ należy do włókna nad $f(a)$, gdyż
$f(a)\in C$. Zatem suma włókien zawiera $A$, a wszystkie są jego
podzbiorami; jest więc równa $A$.
\end{proof}
\begin{przyklad}[Włókna reszty z dzielenia]
Niech $f:\{0,1,2,3,4,5,6,7\}\to\{0,1,2\}$ przypisuje resztę
z dzielenia przez $3$. Włókna to
$\{0,3,6\}$, $\{1,4,7\}$ i $\{2,5\}$.
Każde jest niepuste, żadne dwa nie mają wspólnego elementu i razem
dają wszystkie osiem argumentów. Funkcja jest surjekcją, ale nie
iniekcją, bo na przykład $f(0)=f(3)$.
\end{przyklad}

\begin{figure}[htbp]
\centering
\begin{tikzpicture}[>=Stealth, every node/.style={font=\small}, scale=.98]
 \node[draw,rounded corners,minimum width=1.6cm,minimum height=3.05cm] (dom) at (0,0) {};
 \node[draw,rounded corners,minimum width=1.4cm,minimum height=3.05cm] (cod) at (4.5,0) {};
 \node at (0,1.84) {$A$}; \node at (4.5,1.84) {$B$};
 \node (a1) at (0,.94) {$a_1$}; \node (a2) at (0,.18) {$a_2$};
 \node (a3) at (0,-.58) {$a_3$}; \node (a4) at (0,-1.28) {$a_4$};
 \node (b1) at (4.5,.67) {$b_1$}; \node (b2) at (4.5,-.60) {$b_2$};
 \draw[->,blue!65!black] (a1.east) -- (b1.west);
 \draw[->,blue!65!black] (a2.east) -- (b1.west);
 \draw[->,teal!65!black] (a3.east) -- (b2.west);
 \draw[->,teal!65!black] (a4.east) -- (b2.west);
 \node[align=center,anchor=north] at (2.25,-1.85) {$f^{-1}[\{b_1\}]=\{a_1,a_2\}$,\quad
 $f^{-1}[\{b_2\}]=\{a_3,a_4\}$};
\end{tikzpicture}
\caption{Włókna funkcji rozcinają dziedzinę na klasy argumentów o tej samej wartości.}
\label{fig:fiber-partition-xliv}
\end{figure}

\section{Funkcje stałe, identycznościowe i charakterystyczne}
\begin{definicja}
Dla zbioru $A$ \emph{funkcja identycznościowa} $\id_A:A\to A$ jest określona przez $\id_A(a)=a$. Funkcja $f:A\to B$ jest \emph{stała}, gdy istnieje $b_0\in B$ takie, że $f(a)=b_0$ dla wszystkich $a\in A$. Dla $S\subseteq A$ \emph{funkcja charakterystyczna} $\chi_S:A\to\{0,1\}$ jest określona przez $\chi_S(a)=1$ dla $a\in S$ oraz $\chi_S(a)=0$ dla $a\notin S$.
\end{definicja}
\begin{twierdzenie}[Zbiory i funkcje zero-jedynkowe]\label{thm:characteristic-bijection}
Dla dowolnego zbioru $A$ przyporządkowanie $S\mapsto\chi_S$ jest bijekcją między $\mathcal P(A)$ i $\{0,1\}^A$.
\end{twierdzenie}
\begin{proof}
Każdy podzbiór $S$ wyznacza jednoznaczną funkcję charakterystyczną. Jeśli $\chi_S=\chi_T$, to dla każdego $a\in A$ zachodzi $a\in S\Leftrightarrow\chi_S(a)=1\Leftrightarrow\chi_T(a)=1\Leftrightarrow a\in T$, stąd $S=T$. Jest więc iniekcja. Niech teraz $g:A\to\{0,1\}$. Definiujemy $S:=\{a\in A:g(a)=1\}$. Dla każdego $a$ mamy $\chi_S(a)=g(a)$, bo $g(a)$ może być tylko zerem lub jedynką. Zatem $g=\chi_S$, co dowodzi surjekcji.
\end{proof}
\begin{uwaga}[Przypadek pustej dziedziny]
Dla każdego $B$ istnieje dokładnie jedna funkcja $\varnothing\to B$: relacja pusta. Jeżeli $A\ne\varnothing$, funkcja $A\to\varnothing$ nie istnieje, bo dla dowolnego $a\in A$ nie można wskazać wartości w pustej przeciwdziedzinie. Funkcja $\varnothing\to B$ jest zawsze iniekcją, a surjekcją tylko gdy $B=\varnothing$. Powyższe wynika z definicji, nie z wyjątku od niej.
\end{uwaga}

\begin{uwaga}[Funkcje stałe na pustym zbiorze]
Przy definicji funkcji stałej przyjęliśmy warunek
$\exists b_0\in B\ \forall a\in A\ f(a)=b_0$. Wobec tego
jedyna funkcja $\varnothing\to\varnothing$ \emph{nie jest stała}
w tej konwencji, ponieważ nie można wskazać $b_0$.
Przy spotykanej niekiedy definicji „$f(a)=f(a')$ dla wszystkich
$a,a'\in A$” funkcja ta byłaby stała. Dla niepustego $A$
obie definicje są równoważne. Wyraźne podanie konwencji usuwa
pozorną sprzeczność.
\end{uwaga}

\section{Odwracalność i bijekcje}
\begin{definicja}[Funkcja odwrotna]
Dla $f:A\to B$ funkcję $g:B\to A$ nazywamy \emph{odwrotną} do $f$, jeżeli $g(f(a))=a$ dla każdego $a\in A$ oraz $f(g(b))=b$ dla każdego $b\in B$.
\end{definicja}
\begin{twierdzenie}[Kryterium odwracalności]\label{thm:inverse-iff-bij}
Niech $A,B$ będą zbiorami, a $f:A\to B$ funkcją. Wówczas $f$ ma funkcję odwrotną wtedy i tylko wtedy, gdy jest bijekcją. Wówczas odwrotna jest jednoznaczna i oznaczamy ją $f^{-1}:B\to A$.
\end{twierdzenie}
\begin{proof}
Załóżmy, że $g$ jest odwrotna. Jeśli $f(a)=f(a')$, zastosowanie $g$ daje $g(f(a))=g(f(a'))$, czyli $a=a'$, zatem $f$ jest iniekcją. Dla dowolnego $b\in B$ element $a=g(b)$ spełnia $f(a)=b$, więc $f$ jest surjekcją.

Odwrotnie, niech $f$ będzie bijekcją. Dla każdego $b\in B$ surjekcja zapewnia istnienie $a\in A$ z $f(a)=b$, a iniekcja zapewnia jednoznaczność takiego $a$. Definiujemy $g(b)$ jako ten jedyny element. Wtedy z konstrukcji $f(g(b))=b$, a dla każdego $a$ jednoznaczność przeciwobrazu $f(a)$ daje $g(f(a))=a$. Jeżeli $h$ jest inną odwrotną, to dla $b\in B$ mamy $h(b)=h(f(g(b)))=g(b)$. Zatem $h=g$.
\end{proof}
\begin{uwaga}
Symbol $f^{-1}$ ma dwa zastosowania: przy bijekcji oznacza funkcję odwrotną, a w zapisie $f^{-1}[C]$ dla dowolnej funkcji oznacza przeciwobraz zbioru $C$. Drugi zapis nie wymaga odwracalności. Rozwiniemy to w rozdziale~\ref{ch:function-operations}.
\end{uwaga}


\begin{uwaga}[Notacja złożenia]
Jeżeli $f:A\to B$ i $g:B\to C$, oznaczamy przez $g\circ f:A\to C$
funkcję $a\mapsto g(f(a))$. Formalną definicję i prawa składania
podamy w Rozdziale~\ref{ch:function-operations}. Zapis
$g\circ f=\id_A$ znaczy więc dokładnie, że $g(f(a))=a$
dla wszystkich $a\in A$.
\end{uwaga}

\section{Odwrotność lewostronna i prawostronna}
\begin{definicja}[Jednostronne odwrotności]
Niech $f:A\to B$. Funkcję $g:B\to A$ nazywamy
\emph{lewostronnie odwrotną} do $f$, jeżeli
$g\circ f=\id_A$. Nazywamy ją \emph{prawostronnie odwrotną},
jeżeli $f\circ g=\id_B$. Odwrotność w zwykłym sensie
spełnia oba warunki jednocześnie.
\end{definicja}
\begin{twierdzenie}[Co dają jednostronne odwrotności?]
Niech $A,B$ będą zbiorami, a $f:A\to B$ funkcją.
Jeśli $f$ ma odwrotność lewostronną, to jest iniekcją.
Jeśli ma odwrotność prawostronną, to jest surjekcją.
\end{twierdzenie}
\begin{proof}
Niech $g\circ f=\id_A$ i załóżmy $f(a)=f(a')$.
Po zastosowaniu funkcji $g$ mamy
$g(f(a))=g(f(a'))$, czyli $a=a'$. To jest definicja iniekcji.
Jeśli natomiast $f\circ g=\id_B$, to dla dowolnego $b\in B$
element $a=g(b)$ należy do $A$ i spełnia
$f(a)=f(g(b))=b$. Zatem każdy element $B$ ma przeciwobraz,
co jest definicją surjekcji.
\end{proof}
\begin{uwaga}[Dlaczego nie piszemy od razu odwrotności obu implikacji?]
W teorii mnogości bez aksjomatu wyboru nie każda surjekcja
musi mieć odwrotność prawostronną. Takie odwrócenie wymaga
wybrania po jednym elemencie z każdego włókna funkcji.
Szczegółowo wrócimy do tego problemu przy aksjomacie wyboru.
\end{uwaga}


\begin{twierdzenie}[Dokładne kryterium odwrotności lewostronnej]
\label{thm:left-inverse-converse-xliv}
Niech $A,B$ będą zbiorami, a $f:A\to B$ funkcją.
Istnieje funkcja $g:B\to A$ spełniająca $g\circ f=\id_A$
wtedy i tylko wtedy, gdy $f$ jest iniekcją oraz
$(A\ne\varnothing\ \text{lub}\ B=\varnothing)$.
\end{twierdzenie}
\begin{proof}
Konieczność iniekcyjności wynika z poprzedniego twierdzenia.
Ponadto, jeśli $A=\varnothing$ i $B\ne\varnothing$, nie istnieje
żadna funkcja $B\to A$, co wykazaliśmy w
Twierdzeniu~\ref{thm:function-existence-xliv}. Zatem w tej sytuacji
odwrotność lewostronna nie może istnieć.

Załóżmy teraz, że $f$ jest iniekcją. Jeśli $B=\varnothing$,
warunek istnienia $f:A\to B$ wymusza $A=\varnothing$.
Jedyne funkcje w obie strony mają wtedy puste wykresy,
a ich złożenie jest $\id_{\varnothing}$.
Jeżeli $A\ne\varnothing$, ustalmy $a_0\in A$ i zdefiniujmy
$g:B\to A$ przez
\[
 g(b)=\begin{cases}
 a,&\text{gdy }b=f(a)\text{ dla pewnego }a\in A,\\
 a_0,&\text{gdy }b\notin f[A].
 \end{cases}
\]
Iniekcyjność $f$ zapewnia, że w pierwszym przypadku
$a$ jest jednoznaczny. Oba przypadki się wykluczają i pokrywają
całe $B$. Dla każdego $a\in A$ mamy $g(f(a))=a$,
zatem $g\circ f=\id_A$. Do konstrukcji użyliśmy tylko jednego
ustalonego elementu $a_0$, nie aksjomatu wyboru.
\end{proof}

\begin{przyklad}[Dlaczego pustej dziedziny nie wolno pominąć]
Jedyna funkcja $f:\varnothing\to\{0\}$ jest iniekcją.
Gdyby miała odwrotność lewostronną $g:\{0\}\to\varnothing$,
to musiałaby istnieć wartość $g(0)\in\varnothing$, co jest
niemożliwe. Dla porównania, iniekcja
$f:\{0\}\to\{0,1\}$, $f(0)=0$, ma odwrotność lewostronną
$g(0)=g(1)=0$; brak surjekcyjności nie przeszkadza.
\end{przyklad}

\begin{twierdzenie}[Surjekcje na zbiór skończony mają przekrój]
\label{thm:finite-section-xliv}
Niech $f:A\to B$ będzie surjekcją, a $B$ zbiorem skończonym.
Wówczas istnieje funkcja $g:B\to A$ spełniająca
$f\circ g=\id_B$.
\end{twierdzenie}
\begin{proof}
Jeśli $B=\varnothing$, surjekcja $f:A\to B$ może istnieć
tylko przy $A=\varnothing$; jedyna funkcja pusta $g$
spełnia tezę.
Jeśli $B\ne\varnothing$, wypiszmy jego elementy jako
$B=\{b_1,\dots,b_n\}$, gdzie $n\ge1$.
Surjekcyjność daje dla każdego $i$ niepuste włókno
$f^{-1}[\{b_i\}]$. Wybieramy po jednym elemencie $a_i$
z tych \emph{skończenie wielu} zbiorów: najpierw z pierwszego,
potem z drugiego, aż do $n$-tego. Skończony proces wyboru
można formalnie uzasadnić indukcją po $n$ i nie wymaga aksjomatu
wyboru. Niech $g(b_i)=a_i$. Ponieważ wszystkie $b_i$ są różne,
 $g$ jest dobrze określoną funkcją, a
$(f\circ g)(b_i)=f(a_i)=b_i$ dla każdego $i$.
\end{proof}
\begin{przyklad}[Konkretny przekrój]
Funkcja $f:\{1,2,3,4\}\to\{u,v\}$ dana przez
$f(1)=f(2)=u$, $f(3)=f(4)=v$ jest surjekcją.
Przyporządkowanie $g(u)=2$, $g(v)=4$ jest jej prawostronną
odwrotnością, choć $g$ nie jest odwrotnością dwustronną:
$g(f(1))=2\ne1$.
\end{przyklad}


\begin{twierdzenie}[Przekroje surjekcji a aksjomat wyboru]
\label{thm:surjections-sections-ac-lxvi}
W teorii ZF następujące zdania są równoważne:
\begin{enumerate}[label=(\roman*)]
\item aksjomat wyboru: dla każdego zbioru indeksów $I$ i każdej
rodziny $(X_i)_{i\in I}$ niepustych zbiorów istnieje funkcja $c:I\to\bigcup_{i\in I}X_i$
spełniająca $c(i)\in X_i$ dla wszystkich $i\in I$;
\item dla każdej surjekcji $f:A\to B$ istnieje funkcja
$g:B\to A$ taka, że $f\circ g=\id_B$.
\end{enumerate}
\end{twierdzenie}
\begin{proof}
Załóżmy (i) i niech $f:A\to B$ będzie surjekcją.
Dla każdego $b\in B$ włókno $X_b=f^{-1}[\{b\}]$ jest
niepuste. Aksjomat wyboru zastosowany do rodziny
$(X_b)_{b\in B}$ daje $g(b)\in X_b$. Stąd
$f(g(b))=b$ dla każdego $b\in B$, czyli $f\circ g=\id_B$.
Jeśli $B=\varnothing$, z istnienia funkcji $A\to\varnothing$
wynika $A=\varnothing$ i wybieramy jedyną funkcję pustą.

Załóżmy teraz (ii). Niech $(X_i)_{i\in I}$ będzie rodziną
niepustych zbiorów. Zdefiniujmy rozłączną sumę
\[
 D=\{(i,x):i\in I,\ x\in X_i\}.
\]
Jest ona zbiorem: jej elementy należą do
$I\times\bigcup_{i\in I}X_i$, a wskazany warunek
wyznacza podzbiór przez wyróżnianie. Projekcja
$p:D\to I$, $p(i,x)=i$, jest surjekcją: dla ustalonego
$i\in I$ niepustość $X_i$ daje pewien $x\in X_i$,
a więc $p(i,x)=i$. Nie dokonujemy przy tym jednoczesnego
wyboru wszystkich $x$; dowodzimy tylko zdania
$\forall i\in I\ \exists x\in X_i$.
Z (ii) istnieje $g:I\to D$ z $p\circ g=\id_I$.
Dla każdego $i$ para $g(i)$ ma pierwszą współrzędną $i$,
więc istnieje dokładnie jeden $c(i)\in X_i$ taki, że
$g(i)=(i,c(i))$. Zbiór tych par określa funkcję $c$,
która jest funkcją wyboru. Dla $I=\varnothing$ obie
konstrukcje są puste i argument nadal działa.
\end{proof}
\begin{przyklad}[Dlaczego rozłączna suma jest potrzebna]
Niech $I=\{0,1\}$, $X_0=\{a,b\}$ i $X_1=\{b,c\}$.
Zwykła suma $X_0\cup X_1=\{a,b,c\}$ nie pamięta,
z którego zbioru pochodzi element $b$. Rozłączna suma
$D=\{(0,a),(0,b),(1,b),(1,c)\}$ zachowuje indeksy.
Przekrój $g(0)=(0,a)$, $g(1)=(1,c)$ projekcji
$p:D\to I$ wyznacza funkcję wyboru $c(0)=a$, $c(1)=c$.
Dla dowolnej rodziny niepustych zbiorów istnienie takiego
przekroju jest dokładnie treścią aksjomatu wyboru, a nie
wnioskiem z samej definicji surjekcji w ZF.
\end{przyklad}

\section{Funkcje jako szczególne relacje}
\begin{twierdzenie}[Kryterium wykresu]
Niech $A,B$ będą zbiorami, a $R\subseteq A\times B$ relacją. Relacja $R$ jest wykresem funkcji $A\to B$ wtedy i tylko wtedy, gdy dla każdego $a\in A$ zbiór $\{b\in B:(a,b)\in R\}$ jest jednoelementowy.
\end{twierdzenie}
\begin{proof}
Jeśli $R$ jest wykresem funkcji, to z definicji dla każdego $a$ istnieje dokładnie jedno $b$ z $(a,b)\in R$, czyli wskazany zbiór ma dokładnie jeden element. W drugą stronę jednoelementowość oznacza zarazem istnienie i jednoznaczność wartości dla każdego argumentu, a to jest definicja funkcji.
\end{proof}
\begin{przyklad}[Pierwiastek jako relacja]
Relacja $R=\{(x,y)\in\R^2:y^2=x\}$ nie jest wykresem funkcji $\R\to\R$: dla $x<0$ nie ma żadnego $y$, a dla $x>0$ istnieją dwa rozwiązania. Po ograniczeniu do $A=[0,\infty)$ i $B=[0,\infty)$ relacja jest wykresem funkcji pierwiastkowej: dla każdego $x\ge0$ istnieje dokładnie jeden nieujemny $y$ z $y^2=x$.
\end{przyklad}